\documentclass[11pt,a4paper]{article}

\usepackage[T1]{fontenc}
\usepackage[utf8]{inputenc}
\usepackage{lmodern}
\usepackage{microtype}
\usepackage{graphicx}
\usepackage{booktabs}
\usepackage{hyperref}
\usepackage{url}
\usepackage{xcolor}
\usepackage{enumitem}

\hypersetup{
  colorlinks=true,
  linkcolor=blue,
  citecolor=blue,
  urlcolor=blue
}

\title{%
  \textbf{Incognito:}\\[0.3em]
  \large A Privacy-First, Cross-Platform Tool for\\
  Reviewing and Anonymizing Qualitative Texts\\
  in the Social Sciences
}

\author{%
  Xiaoou Wang\\
  \textit{Maison des Sciences de l'Homme Sud-Est}\\
  \textit{Universit\'{e} C\^{o}te d'Azur}\\
  \texttt{xiaoou.wang@univ-cotedazur.fr}
}

\date{June 2026}

\begin{document}

\begin{center}
  \includegraphics[width=2.2cm]{logo.png}
\end{center}
\vspace{-0.75em}

\maketitle

\begin{abstract}
Qualitative researchers in the social sciences and digital humanities routinely handle
sensitive interview transcripts, field notes, and open-ended survey responses.
Cloud-based anonymization services raise legitimate concerns about data sovereignty,
reproducibility, and compliance with ethics protocols.
We present \textbf{Incognito}, a cross-platform application
developed at the Maison des Sciences de l'Homme Sud-Est (MSHS Sud-Est) that supports
a \emph{human-in-the-loop} workflow for detecting and replacing named entities in
French text entirely on the user's machine or in the browser tab.
The \textbf{desktop edition} combines a React/Electron front end with a persistent local Python NER
service, offering three detection backends---spaCy (\texttt{fr\_core\_news\_sm},
\texttt{fr\_core\_news\_lg}) and CamemBERT NER---plus rule-based detectors for
emails, URLs, phone numbers, and dates.
A complementary \textbf{web interface} (\texttt{web\_interface/}) runs the same review
workflow client-side via Transformers.js and ONNX models (CamemBERT, English BERT, or
custom Hugging Face checkpoints), deployed on GitHub Pages for installation-free access.
Researchers review highlighted spans interactively, toggle replacement at the
category or entity-value level, and export anonymized text with stable placeholders
(e.g.\ \texttt{[PERSON\_1]}), markdown audit reports, and Label Studio
pre-annotations.
A batch mode loads folders of \texttt{.txt} files, supports per-document review
with optional automatic NER, and writes outputs live alongside source files.
Standalone installers for macOS, Windows, and Linux bundle the NER engine via
PyInstaller, requiring no separate Python installation for end users.
We describe the system architecture, workflow, and an evaluation sketch on
French interview-style materials, and discuss limitations regarding indirect
identifiers and model errors.
Source code, release binaries, and citation metadata are publicly available.
\end{abstract}

\noindent\textbf{Keywords:} anonymization, digital humanities, qualitative research,
named entity recognition, French NLP, local-first software, web application, research ethics

\section{Introduction}
\label{sec:introduction}

Anonymization is a recurring obligation in qualitative and mixed-methods research.
Interview excerpts, ethnographic notes, and participant-generated text often contain
direct identifiers (personal names, organizations, precise locations) as well as
quasi-identifiers (rare job titles, small communities, distinctive events) that can
enable re-identification when combined with external knowledge~\cite{elger2014broad}.
Research teams therefore need practical tools that (i)~assist detection of sensitive
strings, (ii)~preserve researcher control over what is removed or generalized,
(iii)~document decisions for ethics committees and data archives, and (iv)~keep
primary data on institutional or personal hardware rather than third-party servers.

General-purpose word processors and manual redaction remain common but scale poorly
to large corpora and offer weak provenance trails.
Meanwhile, high-quality French NER models are now available off the shelf through
libraries such as spaCy~\cite{honnibal2020spacy} and transformer pipelines derived
from CamemBERT~\cite{martin2020camembert,lejeune2020camembertner}.
However, command-line NER outputs are insufficient for humanities workflows:
scholars need visual review, correction of false positives, recovery of false
negatives, category-level policies (e.g.\ replace people but keep organization
mentions), and batch processing aligned with project folder structures.

\textbf{Incognito} addresses this gap as a research prototype and
distribution-ready application (desktop installers and browser deployment) maintained in the context of research
engineering at MSHS Sud-Est\footnote{\url{https://mshs.univ-cotedazur.fr/}}, the
regional Maison des Sciences de l'Homme hosted by Universit\'{e} C\^{o}te d'Azur.
The tool targets French-language qualitative corpora but is designed so that entity
categories and export formats support integration with annotation platforms such as
Label Studio~\cite{tkachenko2020labelstudio} and archival reporting requirements.

Our design principles are:
\begin{itemize}[noitemsep]
  \item \textbf{Local-first processing:} no telemetry and no default calls to
        external NLP APIs; text is sent only to a co-located NER subprocess.
  \item \textbf{Human-in-the-loop control:} automatic suggestions, never blind
        replacement.
  \item \textbf{Traceability:} audit reports listing category decisions and exact
        values replaced.
  \item \textbf{Cross-platform deployment:} the same workflow on macOS, Windows,
        and Linux through Electron installers, or in any modern browser via the
        web interface.
  \item \textbf{Reproducible packaging:} explicit build profiles separating a
        full bundle (spaCy + CamemBERT libraries) from a lighter spaCy-only
        distribution.
\end{itemize}

This paper is a descriptive software report intended for arXiv and citation in
digital humanities methods sections.
It does not claim formal anonymity guarantees; researchers must still review outputs
for indirect identifiers and contextual leaks.

\section{Tool}
\label{sec:tool}

\subsection{Architecture}

Incognito follows a two-process architecture for the desktop edition
(Figure~\ref{fig:architecture}).
The \emph{renderer} is an Electron\footnote{\url{https://www.electronjs.org/}} shell
hosting a React interface built with Vite.
The \emph{NER service} is a long-lived Python process exposing a line-oriented JSON
protocol over standard input/output.
The Electron main process spawns the service at startup, forwards detection requests
from the UI, and handles file-system operations (folder loading, live export of batch
artifacts).

\begin{figure}[ht]
  \centering
  \fbox{\parbox{0.92\linewidth}{
    \small
    \textbf{Renderer (Electron + React)}\\
    Paste / load text $\rightarrow$ entity review UI $\rightarrow$ anonymized preview\\
    $\updownarrow$ IPC\\
    \textbf{Main process}\\
    spawn NER service $\cdot$ folder dialogs $\cdot$ write batch outputs\\
    $\updownarrow$ stdin/stdout JSON\\
    \textbf{NER service (Python)}\\
    spaCy / CamemBERT / rules $\rightarrow$ entity spans with offsets
  }}
  \caption{High-level architecture. All NLP inference runs in the local Python
  subprocess; the renderer never calls cloud APIs.}
  \label{fig:architecture}
\end{figure}

In \emph{development mode}, the main process launches \texttt{scripts/ner.py} with the
project virtual environment.
In \emph{released builds}, a PyInstaller onedir bundle (\texttt{ner-service}) is
placed in application resources; end users do not install Python or spaCy models
separately.
Installers are produced with \texttt{electron-builder} (DMG/ZIP on macOS, NSIS/ZIP
on Windows, AppImage/DEB on Linux).
Released builds ship with a consistent application identity: logo in the integrated
title bar and window icon, plus in-app branding in the review header.
A GitHub Actions workflow builds full standalone artifacts on
\texttt{macos-latest}, \texttt{windows-latest}, and \texttt{ubuntu-latest}.

\subsection{Web interface}

The \texttt{web\_interface/} package is a React application built with Vite.
NER inference runs in a dedicated Web Worker using Transformers.js and ONNX Runtime
WASM; documents never leave the browser except for a one-time download of model
weights from Hugging Face.
It reuses the interactive review, audit report, Label Studio export, and batch
folder workflow (ZIP download instead of live disk writes).
The static build is published to GitHub Pages
(\url{https://xiaoouwang.github.io/Incognito/}) via a dedicated CI workflow.

\subsection{NER backends and entity categories}

The NER service supports three model backends selectable from the UI:
\begin{itemize}[noitemsep]
  \item \textbf{spaCy small} (\texttt{fr\_core\_news\_sm}): faster, lower memory.
  \item \textbf{spaCy large} (\texttt{fr\_core\_news\_lg}): default; richer annotations.
  \item \textbf{CamemBERT NER} (\texttt{Jean-Baptiste/camembert-ner}): transformer
        pipeline with generally higher recall on person and location spans in
        informal French, at higher compute and bundle size.
\end{itemize}

Model labels are mapped to a unified application schema:
\texttt{person}, \texttt{location}, \texttt{organization}, \texttt{misc}, and
\texttt{date} (spaCy).
Complementary \textbf{rule-based detectors} identify email addresses, URLs, phone
numbers, and date expressions (including French month names), which are merged with
model predictions; overlapping spans are resolved with rule spans taking precedence
where configured.

A post-processing step \textbf{splits entity spans at newline boundaries} so that
paragraph breaks are never absorbed into a single replacement---a practical issue we
observed when spaCy grouped tokens across blank lines in interview transcripts.

\subsection{Replacement policy}

Anonymization uses \textbf{stable category placeholders} rather than synthetic
surface forms.
Each unique value within a category receives a consistent tag for the document
(e.g.\ \texttt{[PERSON\_1]}, \texttt{[LOCATION\_2]}).
Replacements are applied by character offset on the original string, not by global
find-and-replace, which preserves unrelated token collisions.

The UI exposes two levels of control:
\begin{enumerate}[noitemsep]
  \item \textbf{Category toggles} --- enable or disable replacement for entire
        classes (people, locations, \ldots).
  \item \textbf{Per-entity-value toggles} --- clicking a chip in the category panel
        disables all occurrences of that value while leaving other entities in the
        category active.
\end{enumerate}

Researchers may also \textbf{manually add or remove spans} in the highlighted view
(e.g.\ to mark a nickname missed by NER or delete a false positive).

\subsection{Exports and interoperability}

Beyond on-screen preview, the tool generates:
\begin{itemize}[noitemsep]
  \item \textbf{Anonymized text} (\texttt{*-anonymized.txt}).
  \item \textbf{Audit report} (\texttt{*-report.md}) with document summary, per-category
        statistics, and exact value$\rightarrow$placeholder mappings.
  \item \textbf{Label Studio JSON} (\texttt{*-label-studio.json}) with
        \texttt{predictions} for span labeling, plus a shared XML labeling config.
\end{itemize}

Command-line utilities in the repository support batch anonymization from Label
Studio exports and folder-scale processing for headless workflows; the desktop
and web applications are the primary interactive paths developed for MSHS Sud-Est researchers.

\section{Workflow}
\label{sec:workflow}

\subsection{Single-document review}

The default workflow proceeds in four stages mirrored in the interface layout:
\begin{enumerate}
  \item \textbf{Source text} --- paste content or edit loaded text directly.
  \item \textbf{Category review} --- run NER, inspect detected categories, toggle
        replacement policies, and click entity chips to include or exclude specific
        values (occurrence counts are shown per value).
  \item \textbf{Highlighted entities} --- read the source with color-coded spans;
        add spans by selection or remove spans by click.
  \item \textbf{Anonymized preview} --- live rendering of placeholder substitutions
        restricted to active categories and entity values.
\end{enumerate}

An audit report modal and clipboard/download actions support ethics documentation.
Optional export to Label Studio creates a pre-annotated task for teams that refine
spans collaboratively before final anonymization.

\subsection{Batch folder mode}

For multi-document projects, \textbf{batch mode} loads a directory of
\texttt{.txt}/\texttt{.text} files (prior output artifacts are filtered automatically).
Each file is opened in the same review UI.
Navigation controls (\emph{previous}/\emph{next file}) preserve per-file entity
states, selected categories, and manual edits in session memory.

When \textbf{auto-run default model} is enabled, NER executes automatically on
each newly opened file that has not yet been processed, using the selected backend
(spaCy large by default).

Critically for archival workflows, batch mode \textbf{writes outputs live} to the
source folder as the researcher works:
anonymized text, audit report, and Label Studio JSON are updated on change
(debounced), so partial progress survives interruption and can be version-controlled
alongside originals.

\subsection{Operational notes}

The first NER invocation after launch may take tens of seconds while models load;
subsequent requests reuse the warm service.
CamemBERT model \emph{weights} ($\sim$400\,MB) are downloaded from Hugging Face on
first use in full builds unless already cached; spaCy pipelines are bundled in
standalone installers.
Researchers should treat audit reports as sensitive documents because they quote
source strings selected for replacement.

\section{Evaluation Sketch}
\label{sec:evaluation}

We report an illustrative evaluation rather than a benchmark with gold-standard
adjudication, which remains future work with ethics-approved corpus releases.

\subsection{Materials}

We used French interview-style excerpts representative of MSHS Sud-Est projects:
semi-structured interviews mentioning persons, institutions, geographic references,
dates, and domain-specific entities (e.g.\ certification bodies, experimental sites).
Internal test sets include multi-paragraph transcripts with orthographic variation
and elided speech typical of oral history transcriptions.

\subsection{Backend comparison (qualitative)}

On the same documents, we compared spaCy small, spaCy large, and CamemBERT:
\begin{itemize}
  \item \textbf{spaCy small} offers the fastest turnaround and smallest memory
        footprint, with occasional missed multi-token organization names.
  \item \textbf{spaCy large} improves coverage of nested entities and is the
        recommended default for interactive review.
  \item \textbf{CamemBERT} tends to propose more person and location spans in
        conversational French; it may over-segment or label misc entities differently,
        motivating the interactive correction UI rather than fully automatic
        anonymization.
\end{itemize}

Rule detectors consistently captured structured identifiers (emails, URLs, phone
patterns) that models sometimes split or skip.

\subsection{Workflow metrics}

For batch sessions, we recorded:
\begin{itemize}[noitemsep]
  \item time-to-first-preview (including model warm-up),
  \item number of manual span additions/removals per 1{,}000 characters,
  \item number of per-entity-value toggles (researcher overrides),
  \item output parity checks ensuring paragraph boundaries are preserved after
        newline-aware span splitting.
\end{itemize}

In internal trials, manual corrections concentrated on \texttt{misc} labels and
rare toponyms; category-level toggles reduced repetitive clicking when entire
classes were out of scope for a given ethics amendment.

\subsection{Limitations and planned work}

The tool does not detect indirect identification (unique life stories, kinship
graphs, profession + city combinations) and does not perform coreference beyond
what each NER model implicitly provides.
We plan:
\begin{itemize}[noitemsep]
  \item a structured user study with MSHS Sud-Est partner laboratories,
  \item adjudicated annotation of a small public French qualitative sample,
  \item optional pseudonymization strategies beyond placeholders,
  \item packaging profiles that pre-download CamemBERT weights for fully offline
        first launch.
\end{itemize}

\section{Availability}
\label{sec:availability}

\paragraph{Software.}
Source code, issue tracker, and citation metadata:
\url{https://github.com/xiaoouwang/Incognito}

\paragraph{Web interface.}
Browser deployment (no install):
\url{https://xiaoouwang.github.io/Incognito/}
Source in \texttt{web\_interface/}; built and published via
\texttt{.github/workflows/deploy-web.yml}.

\paragraph{Releases.}
Standalone \textbf{desktop} binaries (macOS DMG/ZIP for Apple Silicon and Intel, Windows installer,
Linux AppImage/DEB) are distributed via GitHub Releases.
Build instructions support full (\texttt{npm run build:ner:full}) and spaCy-only
(\texttt{npm run build:ner}) NER bundles; metadata in
\texttt{python-dist/ner-service/.bundle-meta.json} records whether CamemBERT
libraries are included.

\paragraph{Citation.}
Please cite this report and the software version used, e.g.:
\begin{quote}
\small
Wang, X. (2026).
\textit{Incognito: A privacy-first, cross-platform tool for reviewing and anonymizing qualitative texts in the social sciences}
(Version 0.2.0) [Computer software].
Maison des Sciences de l'Homme Sud-Est, Universit\'{e} C\^{o}te d'Azur.
\url{https://github.com/xiaoouwang/Incognito}
\end{quote}
\noindent Bib\TeX{} key: \texttt{wang2026incognito}.
A \texttt{CITATION.cff} file is provided for GitHub's citation widget.

\paragraph{License and ethics.}
The application is a research assistant; deployers remain responsible for compliance
with GDPR, institutional review boards, and participant consent.
Contact: \texttt{xiaoou.wang@univ-cotedazur.fr}.

\section*{Acknowledgments}

Developed at the Maison des Sciences de l'Homme Sud-Est (Universit\'{e} C\^{o}te
d'Azur) in support of digital humanities training and qualitative data management.
We thank partner research teams for feedback on batch workflows and audit reporting.

\bibliographystyle{plain}
\bibliography{references}

\end{document}
